% !TEX program = lualatex
% Generated by new_lecture_note.sh
\documentclass[11pt]{article}

\usepackage{fontspec}
\usepackage{microtype}
\usepackage{mathtools,amssymb,amsfonts,amsthm,bm}
\usepackage{enumitem}
\usepackage{booktabs,array,xcolor}
\usepackage{geometry,fancyhdr,setspace}
\usepackage[hidelinks]{hyperref}
\usepackage[nameinlink,noabbrev]{cleveref}

\geometry{margin=1in,headheight=15pt}
\setstretch{1.12}
\setlength{\parindent}{0pt}
\setlength{\parskip}{0.55em plus 0.12em minus 0.08em}
\setlength{\abovedisplayskip}{0.8em plus 0.2em minus 0.1em}
\setlength{\belowdisplayskip}{0.8em plus 0.2em minus 0.1em}
\allowdisplaybreaks[3]
\setlist{leftmargin=*,topsep=0.35em,itemsep=0.15em,parsep=0pt}

% ---------- Metadata ----------
\newcommand{\studentname}{Keshav Ramamurthy}
\newcommand{\coursename}{Stat 150}
\newcommand{\courselongname}{Stochastic Processes}
\newcommand{\lecturenumber}{04}
\newcommand{\lecturetitle}{Lecture 04}
\newcommand{\lecturedate}{\today}

% ---------- Reusable notation ----------
\newcommand{\N}{\mathbb{N}}
\newcommand{\Z}{\mathbb{Z}}
\newcommand{\Q}{\mathbb{Q}}
\newcommand{\R}{\mathbb{R}}
\newcommand{\C}{\mathbb{C}}
\newcommand{\F}{\mathbb{F}}
\newcommand{\K}{\mathbb{K}}
\newcommand{\E}{\mathbb{E}}
\newcommand{\Prob}{\mathbb{P}}
\newcommand{\1}{\mathbf{1}}
\newcommand{\eps}{\varepsilon}
\newcommand{\dd}{\,\mathrm{d}}
\newcommand{\abs}[1]{\left\lvert#1\right\rvert}
\newcommand{\norm}[1]{\left\lVert#1\right\rVert}
\newcommand{\ip}[2]{\left\langle#1,#2\right\rangle}
\newcommand{\set}[1]{\left\{#1\right\}}
\newcommand{\given}{\;\big\vert\;}
\newcommand{\indep}{\perp\!\!\!\perp}
\newcommand{\lto}{\longrightarrow}
\newcommand{\deriv}[2]{\frac{\mathrm{d}#1}{\mathrm{d}#2}}
\newcommand{\pderiv}[2]{\frac{\partial#1}{\partial#2}}
\DeclareMathOperator{\Aut}{Aut}
\DeclareMathOperator{\Hom}{Hom}
\DeclareMathOperator{\Ker}{ker}
\DeclareMathOperator{\im}{im}
\DeclareMathOperator{\rank}{rank}
\DeclareMathOperator{\nullity}{nullity}
\DeclareMathOperator{\Span}{span}
\DeclareMathOperator{\Var}{Var}
\DeclareMathOperator{\Cov}{Cov}

% ---------- Note environments ----------
\newtheoremstyle{notestyle}{0.7em}{0.7em}{\itshape}{}{}{\bfseries}{.5em}{}
\theoremstyle{notestyle}
\newtheorem{theorem}{Theorem}[section]
\newtheorem{lemma}[theorem]{Lemma}
\newtheorem{proposition}[theorem]{Proposition}
\newtheorem{corollary}[theorem]{Corollary}
\theoremstyle{definition}
\newtheorem{definition}[theorem]{Definition}
\newtheorem{example}[theorem]{Example}
\newtheorem{remark}[theorem]{Remark}
\newenvironment{claim}{\par\noindent\textbf{Claim.}\ }{\par}
\newenvironment{idea}{\par\noindent\textit{Idea.}\ }{\par}
\newenvironment{proofsketch}{\par\noindent\textbf{Proof sketch.}\ }{\par}

% ---------- Header ----------
\pagestyle{fancy}
\fancyhf{}
\fancyhead[L]{\coursename}
\fancyhead[C]{\lecturetitle}
\fancyhead[R]{\studentname}
\fancyfoot[C]{\thepage}

\title{\vspace{-1.5em}\textbf{\coursename\ -- \lecturetitle}}
\author{\studentname}
\date{\lecturedate}

\begin{document}
\maketitle

\section{Main idea}
Poisson r.v. $X \sim \text{Poisson}(\lambda)$, with probability generating function
$G_X(s)$ for $s \in [0, 1]$:
$$G_X(s) = \sum_{i=0}^{\infty} \frac{e^{-\lambda}\lambda^i}{i!}s^i =
e^{-\lambda}\sum_{i=0}^{\infty} \frac{(\lambda s)^i}{i!} = e^{-\lambda}e^{\lambda s} =
e^{\lambda(s - 1)} = e^{-\lambda(1 - s)}$$
\section{Galton--Watson branching process (Pinsky \& Karlin \S 3.8.1)}
\par Let $X_n$ be the number of individuals of a population.
\par We often assume $X_0 = 1$ (n is generation n)
Given $X_n = i$  and labelling by $n_1, n_2, n_3, \cdots, n_i$
Suppose each of these individuals gives birth to $\zeta_{n1},
\zeta_{n2},\cdots,\zeta_{ni}$
Then,
$$X_{n+1} = \sum_{j=1}^{X_n} \zeta_{nj}$$
Assume all $\zeta_{nl}$, $n, l \in \Z^{+}$, are independent and have the same distribution.
\par \textbf{Generating function.} Let $G(s) = \E[s^{\zeta}]$ be the offspring PGF
($s \in [0,1]$), with mean $\mu = G'(1)$ and variance $\sigma^2$, and let
$G_n(s) = \E[s^{X_n}]$ (so $G_0(s) = s$, since $X_0 = 1$). A sum of $m$ i.i.d.\ copies
of $\zeta$ has PGF $G(s)^m$, so by the tower property,
$$G_{n+1}(s) = \E\big[\E[s^{X_{n+1}} \given X_n]\big]
= \E\big[G(s)^{X_n}\big] = G_n(G(s)).$$
Each generation just composes $G$ one more time: $G_n = G \circ \cdots \circ G$ ($n$
times). Differentiating at $s = 1$ (where $G(1) = 1$, $G'(1) = \mu$) gives
$\E[X_n] = \mu^n$.

\par \textbf{Two identities.} Everything above runs on two facts.
\par 1. \textit{Tower property}: $\E[\E[X \given Y]] = \E[X]$ --- condition on $Y$, then
average over $Y$, and you are back at the plain average:
$$\E[\E[X \given Y]] = \sum_y \E[X \given Y{=}y]\, \Prob(Y{=}y)
= \sum_x x \sum_y \Prob(X{=}x,\ Y{=}y) = \E[X].$$
\par 2. \textit{Eve's law}: $\Var(Z) = \E[\Var(Z \given Y)] + \Var(\E[Z \given Y])$.
Condition the identity $\E[Z^2] = \Var(Z) + (\E Z)^2$ on $Y$:
$$\E[Z^2 \given Y] = \Var(Z \given Y) + (\E[Z \given Y])^2,$$
average over $Y$, and subtract $(\E Z)^2 = \big(\E[\E[Z \given Y]]\big)^2$; the last two
terms are then exactly the variance of the random variable $\E[Z \given Y]$:
$$\Var(Z) = \E[\Var(Z \given Y)] + \Var(\E[Z \given Y]).$$

\par \textbf{Variance of $X_{n+1}$.} Given $X_n$, the next generation is a sum of
$X_n$ i.i.d.\ copies of $\zeta$, so linearity and independence give
$$\E[X_{n+1} \given X_n] = \mu X_n, \qquad \Var(X_{n+1} \given X_n) = \sigma^2 X_n,$$
and Eve's law turns these into a recursion:
$$\Var(X_{n+1}) = \E[\sigma^2 X_n] + \Var(\mu X_n)
= \sigma^2\, \E[X_n] + \mu^2\, \Var(X_n).$$
When $\mu = 1$ this gives $\Var(X_n) = n\sigma^2$; in general
$\Var(X_n) = \sigma^2\, \mu^{n-1} \frac{\mu^n - 1}{\mu - 1}$ for $\mu \ne 1$.

\section{Reference: notation and facts to remember}
\par \textbf{Notation.}
\begin{itemize}
    \item $X \sim \text{Poisson}(\lambda)$, $\text{Bin}(n,p)$, $\text{Bern}(p)$,
    $\text{Geom}(p)$, $\text{Exp}(\lambda)$: ``$X$ has this distribution.''
    \item i.i.d.: independent and identically distributed.
    \item $X \indep Y$: independent. $\Prob(A \given B) = \Prob(A, B)/\Prob(B)$,
    where ``$,$'' means ``and''.
    \item $\1_A$ (indicator): $1$ if $A$ happens, $0$ otherwise; note
    $\E[\1_A] = \Prob(A)$.
    \item $G_X(s) = \E[s^X]$: probability generating function, $s \in [0,1]$.
    \item $\E[X \given Y{=}y]$ is a \emph{number} (the average of $X$ over outcomes
    with $Y = y$); $\E[X \given Y]$ is a \emph{random variable}, a function of $Y$.
    \item This unit: $X_n$ = size of generation $n$, $\zeta$ = offspring count,
    $\mu = \E[\zeta]$, $\sigma^2 = \Var(\zeta)$, $G$ = offspring PGF.
\end{itemize}
\par \textbf{Expectation.}
\begin{itemize}
    \item Linearity --- no independence needed:
    $\E[aX + bY] = a\,\E[X] + b\,\E[Y]$; sums always split, even with a random number
    of (independent) terms: $\E[X_1 + \cdots + X_N] = \E[N]\,\E[X]$.
    \item $\E[c] = c$; $X \ge 0 \implies \E[X] \ge 0$.
    \item $X \indep Y \implies \E[XY] = \E[X]\,\E[Y]$ (the converse is false).
    \item Tower property: $\E[\E[X \given Y]] = \E[X]$.
\end{itemize}
\par \textbf{Variance and covariance.}
\begin{itemize}
    \item $\Var(X) = \E[X^2] - (\E X)^2 = \E[(X - \E X)^2] \ge 0$.
    \item $\Var(c) = 0$; $\Var(aX) = a^2\,\Var(X)$ --- note the square.
    \item $\Cov(X,Y) = \E[XY] - \E[X]\E[Y]$, symmetric, bilinear,
    $\Cov(X,X) = \Var(X)$.
    \item $\Var(X + Y) = \Var(X) + \Var(Y) + 2\,\Cov(X,Y)$.
    \item Independent $\implies \Cov = 0 \implies$ variances add; $n$ i.i.d.\ terms:
    $\Var(X_1 + \cdots + X_n) = n\sigma^2$.
    \item Eve's law: $\Var(Z) = \E[\Var(Z \given Y)] + \Var(\E[Z \given Y])$.
\end{itemize}
\par \textbf{Generating function facts.}
\begin{itemize}
    \item $G(1) = 1$; $G(0) = \Prob(X = 0)$; $G'(1) = \E[X]$;
    $\Var(X) = G''(1) + G'(1) - \big(G'(1)\big)^2$.
    \item Independent sums multiply: $G_{X+Y} = G_X \cdot G_Y$; $n$ i.i.d.\ copies
    give $G(s)^n$.
    \item Random number of terms: $N \indep$ the $X_i$ gives
    $G_{X_1 + \cdots + X_N}(s) = G_N\big(G_X(s)\big)$ --- this \emph{is} the
    Galton--Watson recursion $G_{n+1}(s) = G_n(G(s))$.
    \item Memorize: $\text{Bern}(p)$: $1 - p + ps$; $\text{Bin}(n,p)$:
    $(1 - p + ps)^n$; $\text{Geom}(p)$: $\frac{ps}{1 - (1-p)s}$;
    $\text{Poisson}(\lambda)$: $e^{\lambda(s-1)}$.
\end{itemize}

\section{Claims, theorems, and proof sketches}
\begin{claim}
State the claim in one precise sentence.
\end{claim}

\begin{proofsketch}
Record the key idea, the first useful equation, and the step that still needs checking.
\end{proofsketch}

\section{Examples and calculations}
\begin{example}[Lecture example]
Write the setup, the calculation, and the conclusion.
\end{example}

\section{Questions and follow-up}
\begin{itemize}
  \item What is still unclear?
  \item Which exercise or reading should I revisit?
  \item TODO:
\end{itemize}

\end{document}
